\setcounter{page}{5}
\setcounter{Exercise}{0}

\chapter*{Feuille d'exercices : Les listes}


% on peut accéder par indice avec des indices positifs en partant du début de la liste ou avec des indices négatifs en partant de la fin de la liste.

%     \begin{tabular}{l c c c c c c c}
%         liste & : &["A", & "B", & "C", & "D", & "E", &"F"]\\
%         indice positif & : & 0 & 1 & 2 & 3 & 4 & 5\\
%         indice négatif & : & -6 & -5 & -4 & -3 & -2 & -1\\
%         \\
%     \end{tabular}




\begin{Exercise}[title={Recherche linéaire}]

\begin{lstlisting}[style=pythonstyle, language=Python]
l = [2,1,11,7,2,6]
x = int(input("Donner une valeur"))
for i in range(len(l)) :
    if l[i] == x:
        print("trouvé à l'index ", i)
    print("non trouvé")
\end{lstlisting}

\Question Quel est le type de la variable l ? Justifiez votre réponse. \\
\emptybox{1cm}

\Question Compléter le tableau d'avancement des variables pour un utilisateur qui a saisi 2 \\

\begin{center}
\large
\begin{tabular}{c|c|c|c|}
\textbf{Tour} & \textbf{i} & \textbf{l[i] == x}  \\
\hline
0 &  &  &  \\
1 &  &  &  \\
2 &  &  &  \\
3 &  &  &  \\
4 &  &  &  \\
5 &  &  &  \\
6 &  &  &  \\
\end{tabular}
\end{center}

\Question Quelles sont les valeurs de retours de la fonction recherche pour les variables suivantes : \\
\begin{itemize}
    \item l = [2, 1, 11] et x= 7
    \item l = [] et x = 3
    \item l = [1, 0] et x = 5
    \item l = [6, 1, 11] et x = 11
\end{itemize}

\Question Expliquer l'objectif et le fonctionnement de cette programme en une phrase.\\

\end{Exercise}

\begin{Exercise}[title={Trouver un maximum}]

\begin{lstlisting}[style=pythonstyle, language=Python]
tab = [9, 100, 8, 48, 73]
imax, vmax = 0, tab[0]
for i in range(len(tab)):
    if tab[i] > vmax:
        vmax = tab[i]
        imax = i
print(vmax, imax) 
\end{lstlisting}

\Question Quel type de parcours est utilisé à la ligne 3 ?\\
\Question Que représente la variable imax ? \\
\Question Compléter le tableau d'avancement des variables\\
 \begin{center}
\large
\begin{tabular}{c|c|c|c|c|c|}
\textbf{Tour} & \textbf{i} & \textbf{tab[i]} & \textbf{tab[i] > vmax} & \textbf{vmax} & \textbf{imax } \\
\hline
0 &  &  &  & &\\
1 &  &  &  & &\\
2 &  &  &  & &\\
3 &  &  &  & &\\
4 &  &  &  & &
\end{tabular}
\end{center}

\Question Quelles sont les valeurs affichées pour les lsites suivantes:
\begin{itemize}
 \item l = [10, 4, 8, 88, 10]
 \item l = [-10, -2, -10]
\end{itemize}

\Question Expliquer l'objectif et le fonctionnement de cette programme en une phrase.\\
\end{Exercise}
\begin{Exercise}[title={Vérification tri}]

\begin{lstlisting}[style=pythonstyle, language=Python]
tab = [0, 5, 8, 8, 9]
tri = True
for i in range(len(tab)-1):
    if tab[i] > tab[i+1]:
        tri = False
print(tri)
\end{lstlisting}

\Question Combien y a-t-il de structures de contrôle dans le code ci-dessus ?\\

\Question Quel type de parcours est utilisé ?\\
\Question Compléter le tableau d'avancement des variables

\begin{center}
\large
\begin{tabular}{c|c|c|c|c|}
\textbf{Tour} & \textbf{i} & \textbf{tab[i]} & \textbf{tab[i+1]} & \textbf{tab[i] > tab[i+1]}  \\
\hline
0 &  &  &  & \\
1 &  &  &  & \\
2 &  &  &  & \\
3 &  &  &  & \\
4 &  &  &  & \\
5 &  &  &  & \\
\end{tabular}
\end{center}

\Question Quelles sont les affichages pour les tab suivants : \\
\begin{itemize}
    \item tab = [8, 12, 4]
    \item tab = [-1, 4]
    \item tab = [5]
\end{itemize}

\Question Expliquer l'objectif et le fonctionnement de cette programme en une phrase.\\
\end{Exercise}
% \begin{Exercise}[title={Recherche linéaire}]

% \begin{lstlisting}[style=pythonstyle, language=Python]
% def recherche(l, x):
%     for i in range(len(l)) :
%         if l[i] == x:
%             return l[i]
%     return -1
% \end{lstlisting}

% \Question Quel est le type du paramètre l ? Justifiez votre réponse. \\
% \emptybox{1cm}

% \Question Compléter le tableau d'avancement des variables pour un appel de la fonction recherche([2,1,11,7,2, 6]) \\

% \begin{center}
% \large
% \begin{tabular}{c|c|c|c|}
% \textbf{Tour} & \textbf{k} & \textbf{v} & \textbf{k == v}  \\
% \hline
% 0 &  &  &  \\
% 1 &  &  &  \\
% 2 &  &  &  \\
% 3 &  &  &  \\
% 4 &  &  &  \\
% \end{tabular}
% \end{center}

% \Question Quelles sont les valeurs de retours de la fonction recherche pour les appels suivants : \\
% \begin{itemize}
%     \item recherche([2, 1, 11], 7)
%     \item recherche([], 3)
%     \item recherche([1, 0], 5)
%     \item recherche([6, 1, 11], 11)
% \end{itemize}

% \Question Expliquer ce que l'objectif et le fonctionnement de cette fonction en 1 paragraphe.\\

% \end{Exercise}

% \begin{Exercise}[title={Trouver un maximum}]

% \begin{lstlisting}[style=pythonstyle, language=Python]
% def recherche_max(tab):
%     imax, vmax = 0, tab[0]
%     for i in range(len(tab)):
%         if tab[i] > vmax:
%             vmax = tab[i]
%             imax = i
%     return vmax, imax
% \end{lstlisting}

% \Question Quel type de parcours est utilisé à la ligne 3 de la fonction ?\\
% \Question Que représente la variable imax ? \\
% \Question Écrire l'appel de la fonction recherche\_max() qui prendra en paramètre la liste [9, 100, 8, 48, 73]\\
% \Question Compléter le tableau d'avancement des variables pour un appel de la fonction recherche\_max() qui prendra en paramètre la liste [9, 100, 8, 48, 73].\\
%  \begin{center}
% \large
% \begin{tabular}{c|c|c|c|c|c|}
% \textbf{Tour} & \textbf{i} & \textbf{tab[i]} & \textbf{tab[i] > vmax} & \textbf{vmax} & \textbf{imax } \\
% \hline
% 0 &  &  &  & &\\
% 1 &  &  &  & &\\
% 2 &  &  &  & &\\
% 3 &  &  &  & &\\
% 4 &  &  &  & &
% \end{tabular}
% \end{center}

% \Question Quelles sont les valeurs de retours de la fonction recherche\_max pour les appels suivants :
% \begin{itemize}
%  \item recherche\_max([10, 4, 8, 88, 10]):
%  \item recherche\_max([-10, -2, -10]): \\
% \end{itemize}

% \Question Expliquer ce que l'objectif et le fonctionnement de cette fonction en 1 paragraphe.\\
% \end{Exercise}
% \begin{Exercise}[title={Vérification tri}]

% \begin{lstlisting}[style=pythonstyle, language=Python]
% def verifie(tab):
%     for i in range(len(tab)-1):
%         if tab[i] > tab[i+1]:
%             return False
%         return True
% \end{lstlisting}

% \Question Combien y a-t-il de structures de contrôle dans le code ci-dessus ?\\

% \Question Quel type de parcours est utilisé à la ligne 2 de la fonction ?\\
% \Question Compléter le tableau d'avancement des variables pour un appel de la fonction verifie([0, 5, 8, 8, 9]).\\

% \begin{center}
% \large
% \begin{tabular}{c|c|c|c|c|}
% \textbf{Tour} & \textbf{i} & \textbf{tab[i]} & \textbf{tab[i+1]} & \textbf{tab[i] > tab[i+1]}  \\
% \hline
% 0 &  &  &  & \\
% 1 &  &  &  & \\
% 2 &  &  &  & \\
% 3 &  &  &  & \\
% 4 &  &  &  & \\
% 5 &  &  &  & \\
% \end{tabular}
% \end{center}

% \Question Quelles sont les valeurs de retours de la fonction verifie pour les appels suivants : \\
% \begin{itemize}
%     \item verifie([8, 12, 4])
%     \item verifie([-1, 4])
%     \item verifie([5])
% \end{itemize}

% \Question Expliquer ce que l'objectif et le fonctionnement de cette fonction en 1 paragraphe.\\
% \end{Exercise}



